\thispagestyle{empty}
\vspace*{\fill}
\begin{center}
    \textcolor{univblue}{\rule{\textwidth}{3pt}} \\[0.3cm]
    {\LARGE\bfseries Design and Development of a Web and Mobile Application for HMI Stars : Expert Accountant} \\[0.4cm]
    \textcolor{univblue}{\rule{\textwidth}{3pt}}
\end{center}

\vspace{0.5cm}

\begin{otherlanguage}{french}
\begin{spacing}{1.1}
\noindent \textbf{Résumé : } 
Ce projet de fin d'études porte sur la conception et la réalisation d'un écosystème applicatif de communication et de gestion en temps réel nommé \textbf{HMI Stars}. Conçu pour le cabinet de conseil \textbf{HMI Stars Consulting}, cet écosystème vise à simplifier, automatiser et centraliser les échanges administratifs, sociaux et financiers entre le cabinet et ses entreprises clientes. L'architecture repose sur un monorepo Flutter contenant deux applications distinctes connectées à une base de données \textbf{Supabase} : une plateforme web d'administration destinée aux gestionnaires du cabinet, et une application mobile destinée aux gérants des entreprises clientes. Cet écosystème intègre des mécanismes de sécurité via les politiques RLS (Row Level Security) de PostgreSQL et résout les boucles de redirection d'authentification sur le web grâce à la mise en place d'un protocole de routage adaptatif et à l'exploitation du flux d'authentification implicite de Supabase.

\noindent \textbf{Mots-clés : } Flutter, Supabase, PostgreSQL, Temps Réel, Messagerie d'Entreprise, RLS.
\end{spacing}
\end{otherlanguage}

\vspace{0.6cm}

\begin{spacing}{1.1}
\noindent \textbf{Abstract: } 
This graduation project focuses on the design and implementation of a real-time communication and administrative management ecosystem named \textbf{HMI Stars}. Built specifically for the consulting enterprise \textbf{HMI Stars Consulting}, this ecosystem aims to simplify, automate, and centralize social, administrative, and financial exchanges between the enterprise and its corporate clients. The system architecture is established on a Flutter monorepo containing two distinct applications connected to a single \textbf{Supabase} backend: a web-based administration platform for the enterprise's managers, and a cross-platform mobile application for client managers. This ecosystem integrates advanced security mechanisms via PostgreSQL Row Level Security (RLS) policies and ensures authentication flow consistency on the web by implementing an adaptive routing protocol combined with Supabase's implicit auth flow.

\noindent \textbf{Keywords: } Flutter, Supabase, PostgreSQL, Real-time, Enterprise Messaging, RLS.
\end{spacing}

\vspace{0.6cm}

\begin{otherlanguage}{arabic}
\begin{spacing}{1.1}
\noindent \textbf{الخلاصة: } 
يتمحور مشروع نهاية الدراسة هذا حول تصميم وإنجاز منظومة برمجية متكاملة للتواصل والإدارة في الوقت الفعلي تحت اسم \textbf{HMI Stars}. صُممت هذه المنظومة لفائدة مكتب الاستشارات \textbf{HMI Stars Consulting}، وتهدف إلى تبسيط، أتمتة ومركزة التبادلات الإدارية، الاجتماعية والمالية بين المكتب والشركات العميلة له. تعتمد بنية النظام على مستودع أحادي للغة Flutter يحتوي على تطبيقين منفصلين متصلين بقاعدة بيانات \textbf{Supabase}: منصة ويب للإدارة مخصصة لمديري المكتب، وتطبيق جوال مخصص لمديري الشركات العميلة. تدمج هذه المنظومة آليات حماية متطورة عبر سياسات أمن مستوى الصف (RLS) لقاعدة بيانات PostgreSQL وتضمن اتساق تدفقات التحقق من الهوية على الويب من خلال إرساء بروتوكول توجيه متكيف واستغلال تدفق التحقق الضمني لـ Supabase.

\noindent \textbf{الكلمات المفاتيح: } فلاتر، سوبابيس، بوستجريس كيو إل، الوقت الفعلي، مراسلات المؤسسات، أمن مستوى الصف.
\end{spacing}
\end{otherlanguage}
\vspace*{\fill}
